% source: manti.01.ma, 1ère SM, devoir surveillé (Drive 0B6pfSJcEVng0UDRFUEx6N21ycE0). Ex.3-4 géométrie (hors thème)
\ex{Suites}{UDRF}{Ex. 1}{6}
Soit $(U_n)_n$ la suite définie par $U_0=-1$, $U_1=1$ et $3U_{n+2}=5U_{n+1}-2U_n$.
\begin{enumerate}
\item On pose $W_n=U_{n+1}-U_n$. Montrer que $(W_n)_n$ est une suite géométrique de raison $q=\dfrac23$.
\item En déduire que $U_n=5-6\left(\dfrac23\right)^n$.
\end{enumerate}
\fin
\ex{Suites}{UDRF}{Ex. 2}{9}
Soient $(U_n)_n$ et $(V_n)_n$ les suites définies par
$U_0=1,\ U_{n+1}=\dfrac{3U_n+V_n}{4}$ \quad et \quad $V_0=6,\ V_{n+1}=\dfrac{U_n+5V_n}{6}$.
\begin{enumerate}
\item Calculer $U_1$ et $V_1$.
\item On pose $W_n=V_n-U_n$. Montrer que $(W_n)_n$ est une suite géométrique.
\item On pose $T_n=2U_n+3V_n$. Montrer que $(T_n)_n$ est constante et préciser sa valeur.
\item En déduire $U_n$ et $V_n$ en fonction de $n$.
\end{enumerate}
\fin
